\documentclass{article}
\usepackage{amsmath,amssymb,hyperref}
\title{A diagnostic intervention can identify the wrong forecast}
\author{}
\date{26 November 2026}
\begin{document}
\maketitle

\paragraph{Audit.}
This is an exact finite R3b probe, not a physical study, an independent
test or completion of R3a. A site's reserve must be staged before the
next mission; the choice cannot wait to observe that mission's failure.
The reserve offers either of two lawful lanes, with safe slow fallback
if staged for the wrong lane. A calibration law is drawn from the pinned
two-law family and stays fixed across episodes. Episodes are independent
conditional on the law and the selected load setting.
An earlier optional load admission changes both access to a later observation
and potentially the plant law itself. No assumption that these two effects
are separate in real systems is made. Calibration costs are measured in
the same normalized service-loss units as slower fallback. The actuated
geometry and noise probabilities are stipulated.

\paragraph{Constructive observational pair.}
Let $X$ be a fair bit for the initial lane, let $Z$ be the actual
exchange parity and $Y=X\oplus Z$ the final lane, with $Z$ independent
of $X$. In an episode with optional work admitted, the ordinary
diagnostic can observe one of $X,Y$, not both: its setup for either
slot and the ongoing work preclude acquisition at the other slot.
For every $p=\Pr(Z=1)$, either marginal is fair. The slot choice can
depend on all earlier episodes and its own coins. Conditioning on all
of that history, the next single reading is still fair, proving equality
of the entire experiment transcript by induction.
Declining the optional load before the first slot enables joint reads
but changes the law to $q=\Pr(Z=1\mid\text{unloaded})$. In two source
models $q=1/2$ while $p=1/2-a$ or $p=1/2+a$ ($0<a\le1/2$).
The unloaded joint pair has the same full distribution in both models;
choosing it adaptively cannot distinguish them either. For a single
structural construction, take a fresh uniform $U$ and threshold it by
$q$ when unloaded or by the corresponding $p$ when loaded.
The intervention pair is entirely observed but is not the target mission.
The two laws require opposite preparations: placing reserve on lane 1
wins when $p>1/2$, and lane 0 wins when $p<1/2$. A wrong forecast
costs $|2p-1|$ additional expected fallback loss relative to known-law
preparation. Without a new loaded experiment or an assumption linking
regimes, indefinitely more of either historical experiment cannot fix it.

\paragraph{Sharp transport decision, not just a missing parameter.}
When the unloaded law $q$ is known and there is independent warrant
for $|p-q|\le\delta$, the exact identified interval in this family is
$[l,u]=[\max(0,q-\delta),\min(1,q+\delta)]$. Every point is realized
by the explicit pair construction. This is sensitivity to a shift in
the source law, not a statistical confidence interval or a way to learn
$\delta$ from the same uninformative data. A sign is pinned when the
interval lies to one side of $1/2$, even if $p$ remains unknown.
Set $A=(1-2l)_+$ and $B=(2u-1)_+$. If both are positive, staging
on lane 1 with probability $B/(A+B)$ equalizes the worst endpoint
regrets and gives $AB/(A+B)$. The same formula gives zero when only
one is positive; both zero means any choice. The lower bound holds
because any distribution has worst regret at least
$\max(Ax,B(1-x))$ in the crossing case, while this choice attains it.
For a wholly one-sided interval simply pick its known optimal lane.

\paragraph{A loaded experiment with protected time.}
Reserve a distinct diagnostic slot at cost $c>0$ for one more
calibration episode without dropping the original load; it displaces
a different optional job, rather than silently waiving a hard duty.
The slot permits paired transition reads in the loaded regime.
One successful pair produces $Z$ with an independently calibrated
bit flip of probability $\epsilon\le1/2$; a whole pair can fail with
probability $d$, independently of parity. For symmetric alternatives
$p=1/2\pm a$, the observed non-erased bit probabilities under the two
laws are $1/2\pm a(1-2\epsilon)$. The best worst-law classification
error is $1/2-a(1-2\epsilon)(1-d)$; the mirrored likelihood test
achieves it, and the average under equal priors lower-bounds every
maximum. As the future loaded episode is distinct, one such protected
read is worthwhile over none exactly when
$c<2a^2(1-2\epsilon)(1-d)$. The formula does not transport to
state-dependent missingness. In the selected-channel negative control,
a fair reported bit is entirely uninformative conditional on delivery
yet the calibrated presence bit alone distinguishes the two laws.
In a nonsymmetric control the equal-prior likelihood error is strictly
less than the optimal worst-law error.

\paragraph{Adaptive calibration before staging.}
With at most $H$ further protected trials and common per-trial cost,
the symmetric two-law problem has an exact posterior recursion.
At time-to-go $n$ and posterior $z$ on the positive law, stop at loss
$2a\min(z,1-z)$, or pay $c$ and take the expectation of the previous
value at the updated posterior. A rejected sample may be an explicit
erasure and leaves the posterior unchanged in the symmetric case.
With initial posterior $1/2$, symmetric optimal play has equal risks
under either law, so the recursion is minimax. Compare with every
fixed sample count $0,\dots,H$ and its optimal likelihood decoder.
For $H=5$, $a=1/4$, $\epsilon=0$, $d=1/3$ and $c=1/64$,
adaptive stopping yields $8675/62208$ while the best fixed count
costs $209/1296$. A separate unnormalized binary likelihood tree
agrees exactly; there is no advantage over equally informed
controlled sensing. Both include the price of taking more samples.
There is no claim of improving a more general stopping rule.

\paragraph{Relation to prior work and replay.}
The project's earlier parity example lets the costly unloaded scan
observe the target exchange law. This one falsifies that transport step
even though the unloaded experiment has complete support.
\href{https://proceedings.mlr.press/v286/marmarelis25a.html}{Marmarelis
et al.} also bound off-policy outcomes using a sensitivity model, but
their trajectory confounding bound is not the asserted absolute
between-regime shift bound here.
\href{https://arxiv.org/abs/1205.0858}{Nitinawarat, Atia and Veeravalli}
already study controlled and sequential sensing. The point here is the
requirement to perform calibration in the right intervention regime, not
a new generic stopping method.
The exact replay compares 289 shift cases, 1377 symmetric one-shot
cases and 405 finite stopping comparisons. It additionally constructs
a witness for strict stopping value and four explicit asymmetry
controls. The loaded diagnostic, all transport parameters, and
persistent access costs remain hypotheses for a source-valid study.
\end{document}
